% Table: Cost-benefit analysis of expanding economics teacher supply (Chapter 4)
\begin{table}[htbp]
\centering
\caption{Cost-benefit analysis of expanding economics teacher supply}
\label{tab:mvpf}
\footnotesize\setlength{\tabcolsep}{4pt}
\begin{tabular}{l l r l}
\toprule
\multicolumn{4}{c}{\textbf{Panel A: Marginal value of public funds}} \\
\midrule
Parameter    & Description                              & Value         & Source \\
\midrule
$\delta$     & Teacher $\rightarrow$ sustained provision & 1             & Assumed \\
$\pi_1$      & Availability $\rightarrow$ A-level       & 0.057            & Table~\ref{tab:main} \\
$\beta_1$    & A-level $\rightarrow$ gross PV lifetime return  & \pounds 185,000           & Table~\ref{tab:main} Panel~C, adjusted\\
$\beta_1^p$  & A-level $\rightarrow$ net private lifetime return & \pounds 88,000           & Table~\ref{tab:main} Panel~C, adjusted\\
$\beta_1^f$  & A-level $\rightarrow$ exchequer return      & \pounds 97,000           & Table~\ref{tab:main} Panel~C, adjusted\\
$N$          & Cohort size per school                   & 46            & Data \\
$C$          & PV cost per teacher (salary + training)  & \pounds 645,000           & DfE; NFER \\
\addlinespace
             & \textbf{MVPF}                            & Infinite   &  \\
             & \textbf{Fiscal surplus per teacher}      & \pounds 1,540,000   &  \\
\midrule
\multicolumn{4}{c}{\textbf{Panel B: Universal access counterfactual}} \\
\midrule
             &                                          & Current       & Universal access \\
\midrule
             & A-level econ students per cohort         & 15,026            & 19,018 \\
             & Economics undergraduates per cohort      & 5,714            & 6,321 \\
\addlinespace
             & \% female among economics undergraduates & 28.2            & 28.3 \\
             & \% FSM among economics undergraduates    & 6.8            & 6.9 \\
             & \% outside London/SE among econ.\ undergraduates
                                                        & 50.6            & 53.3 \\
\addlinespace
             & Aggregate gross PV earnings gain per cohort &            & \pounds 740 m \\
\bottomrule
\end{tabular}
\begin{minipage}{\textwidth}
\smallskip
\footnotesize
\textit{Notes:} Panel~A reports the components of the MVPF calculation. The teacher-to-provision conversion $\delta$ is assumed to equal one: a hired teacher is taken to translate into sustained provision. Section~\ref{sec:Policy} reports the break-even conversion rate. The lifetime return parameters are the composition-adjusted implied values of Section~\ref{sec:heterogeneity}, which reweight the Table~\ref{tab:main} Panel~C returns to the composition of never-offer schools. The MVPF numerator uses the net private return and the denominator the exchequer return; the gross return is shown for context and approximately equals their sum. The MVPF equals beneficiaries' willingness to pay (private lifetime earnings gain) divided by the net government cost (teacher salary and training costs minus additional tax revenue from higher earnings). An MVPF above 1 indicates the policy is welfare-improving. If the fiscal return exceeds the cost, the denominator is negative and the MVPF is infinite --- the policy pays for itself. Panel~B reports a counterfactual exercise in which all schools are assumed to offer A-level economics. Additional A-level students and economics undergraduates are computed from the first-stage estimates applied to the population of students currently in schools that do not offer economics. Compositional shifts apply each group's take-up effect of availability to its no-access population --- the sex-specific effects (2.8 and 9.2 percentage points, Section~\ref{sec:Policy}) for the sex row and the pooled effect (5.7 points) for the FSM and region rows --- with conversion to economics degrees at the 15.2 percentage point gateway effect. Shares are computed among students with the relevant characteristic recorded; under the sex-specific effects the implied additions are slightly fewer than the pooled totals in the count rows. The aggregate earnings gain is gross; the net individual analogue, applying the ratio of the net to the gross return in Table~\ref{tab:main} Panel~C, is roughly \pounds 350~million. FSM denotes eligibility for free school meals, a standard administrative marker of low socio-economic family background.
\end{minipage}
\end{table}
